\chapter{Legislativa a referenční hodnoty pro IEQ}
\label{app:legislativa}

Tato příloha doplňuje kapitolu \textit{Teoretická východiska} o~podrobné přehledové tabulky. Tabulka~\ref{tab:legislativa_ieq_dokumenty} uvádí relevantní dokumenty ČR, EU a~WHO a~jejich pokrytí jednotlivých domén IEQ. Tabulka~\ref{tab:iaqlimity_cr_eu_who} shrnuje limitní a~doporučené hodnoty pro vybrané látky vnějšího ovzduší, které jsou relevantní pro IAQ.

\begin{table}[p]

\begin{center}
\setlength{\tabcolsep}{3pt}
\renewcommand{\arraystretch}{1.05}
\scriptsize
\caption{Dokumenty relevantní pro IEQ a jejich pokrytí (řádek = dokument)}\label{tab:legislativa_ieq_dokumenty}
\begin{tabularx}{\textwidth}{@{}
>{\raggedright\arraybackslash}p{0.33\textwidth}
  >{\raggedright\arraybackslash}p{0.24\textwidth}
  X@{}}
\toprule
\textbf{Dokument (úroveň, typ)} & \textbf{Prostředí / domény} & \textbf{Vymezené oblasti a typicky sledované veličiny} \\
\midrule

Zákon č.~258/2000~Sb. &
Veřejné zdraví; IAQ, akustika &
Rámcový zákon k ochraně veřejného zdraví (kompetence, postupy); zmocňuje prováděcí předpisy pro hluk a vibrace. \\
Vyhláška č.~43/2025~Sb. &
Indoor (pobytové); IAQ, mikroklima &
T, RH, CO$_2$ (explicitní limity 1\,200/2\,500 ppm), PM$_{2.5}$, chemické ukazatele; nahrazuje zrušenou vyhl.~6/2003~Sb. \\

Zákon č.~283/2021~Sb. &
Stavby; vizuální komfort &
Technické požadavky na stavby (mj. požadavky na denní/umělé osvětlení v kontextu návrhu a užívání staveb; detaily často odkazují na ČSN). \\

Nařízení vlády č.~361/2007~Sb. &
Pracoviště; IAQ, mikroklima &
CO$_2$ (expoziční limit), mikroklima (T, RH, rychlost proudění) aj. dle zátěže. \\

Nařízení vlády č.~272/2011~Sb. &
Indoor/Outdoor; akustika &
Hygienické limity hluku (pracoviště, chráněné vnitřní/venkovní prostory) a vibrací; způsob měření a hodnocení. \\

Metodický návod ZD~25/2013 &
Metodika měření &
Postupy měření a hodnocení (umístění čidel, doba měření, vyhodnocení). \\

Zákon č.~201/2012~Sb. &
Vnější ovzduší; IAQ &
Imisní/cílové/prahové hodnoty (SO$_2$, NO$_2$, CO, benzen, Pb, PM$_{10}$, PM$_{2.5}$, O$_3$). \\

\midrule
Směrnice 2008/50/ES &
Vnější ovzduší; IAQ &
Harmonizace limitů/cílů a pravidel hodnocení/reportingu (PM, NO$_2$, SO$_2$, CO, benzen, Pb, O$_3$). \\

Směrnice 2002/49/ES &
Venkovní prostředí; akustika &
Hodnocení a řízení hluku ve venkovním prostředí (strategické hlukové mapy, akční plány). \\

Směrnice 2010/31/EU &
Budovy; IEQ, vizuální komfort &
Požaduje zohlednění kvality vnitřního prostředí při opatřeních na energetickou náročnost; pracuje s technickými systémy budovy vč. zabudovaného osvětlení. \\

Směrnice (EU) 2024/1275 &
Budovy; IEQ, vizuální komfort &
Aktualizovaná EPBD; zohledňuje technické systémy a automatizaci budov včetně osvětlení. \\

\midrule
WHO Air Quality Guidelines &
Zdravotní doporučení (benchmark) &
Doporučené úrovně expozice pro PM$_{2.5}$/PM$_{10}$, O$_3$, NO$_2$, SO$_2$, CO (různé doby průměrování). \\

WHO indoor: selected pollutants &
Indoor; IAQ &
Referenční hodnoty pro vybrané látky v interiéru (např. formaldehyd, CO, NO$_2$ aj.). \\

WHO indoor: dampness and mould &
Indoor; vlhkost/plísně &
Prevence a hodnocení rizik (důraz na zdroje vlhkosti a kontext budovy). \\

WHO indoor: household fuel combustion &
Indoor; spalování &
Doporučení ke snížení expozice (typicky PM a CO) dle paliva/technologie. \\

\bottomrule
\end{tabularx}
\end{center}

\footnotesize \textit{Zdroj:} \parencite{wolterskluwerZakon2582000,info@wolterskluwer.czVyhlaska432025,info@wolterskluwer.czZakon2832021,wolterskluwerNarizeni36120072007,info@wolterskluwer.czNarizeni2722011,MetodickyNavodZD252013,wolterskluwerZakon20120122012,SmerniceEvropskehoParlamentu2008,SmerniceEvropskehoParlamentu2002,SmerniceEvropskehoParlamentu2010,SmerniceEvropskehoParlamentu2024a,worldhealthorganizationWHOGlobalAir2021,adamkiewiczWHOGuidelinesIndoor2010,heseltine2009guidelines,world2014guidelines}
\end{table}

\begin{table}[p]

\begin{center}
\setlength{\tabcolsep}{3pt}
\renewcommand{\arraystretch}{1.05}
\scriptsize
\caption{Srovnání limitních a doporučených hodnot pro vybrané látky vnějšího ovzduší}\label{tab:iaqlimity_cr_eu_who}
\begin{threeparttable}
\begin{tabularx}{\textwidth}{@{}
>{\raggedright\arraybackslash}p{0.10\textwidth}
  >{\raggedright\arraybackslash}p{0.16\textwidth}
  >{\raggedright\arraybackslash}p{0.09\textwidth}
  >{\raggedright\arraybackslash}p{0.09\textwidth}
  >{\raggedright\arraybackslash}p{0.10\textwidth}
  X@{}}
\toprule
\textbf{Látka} & \textbf{Doba (jedn.)} & \textbf{ČR} & \textbf{EU} & \textbf{WHO} & \textbf{Pozn.} \\
\midrule

SO\textsubscript{2} & 1 h ($\mu\text{g}\,m^{-3}$) & 350 & 350 & — & $\leq$24\,\texttimes{}/rok (ČR/EU) \\
SO\textsubscript{2} & 24 h ($\mu\text{g}\,m^{-3}$) & 125 & 125 & 40 & $\leq$3\,\texttimes{}/rok (ČR/EU) \\
NO\textsubscript{2} & 1 h ($\mu\text{g}\,m^{-3}$) & 200 & 200 & 200 & ČR/EU: $\leq$18\,\texttimes{}/rok; WHO: Tab. 0.2 \\
NO\textsubscript{2} & 24 h ($\mu\text{g}\,m^{-3}$) & — & — & 25 & WHO: Tab. 0.1 \\
NO\textsubscript{2} & 1 rok ($\mu\text{g}\,m^{-3}$) & 40 & 40 & 10 & — \\

CO & max. denní 8 h ($\text{mg}\,m^{-3}$) & 10 & 10 & 10 & ČR/EU: klouzavý 8h; WHO: Tab. 0.2 \\
CO & 24 h ($\text{mg}\,m^{-3}$) & — & — & 4 & WHO: Tab. 0.1 \\

Benzen & 1 rok ($\mu\text{g}\,m^{-3}$) & 5 & 5 & — & — \\
Pb & 1 rok ($\mu\text{g}\,m^{-3}$) & 0{.}5 & 0{.}5 & — & — \\

PM\textsubscript{10} & 24 h ($\mu\text{g}\,m^{-3}$) & 50 & 50 & 45 & ČR/EU: $\leq$35\,\texttimes{}/rok \\
PM\textsubscript{10} & 1 rok ($\mu\text{g}\,m^{-3}$) & 40 & 40 & 15 & — \\

PM\textsubscript{2.5} & 24 h ($\mu\text{g}\,m^{-3}$) & — & — & 15 & WHO: Tab. 0.1 \\
PM\textsubscript{2.5} & 1 rok ($\mu\text{g}\,m^{-3}$) & 20 & 25 & 5 & EU: 25 (2015); ind.: 20 (2020) \\

O\textsubscript{3} & max. denní 8 h ($\mu\text{g}\,m^{-3}$) & 120 & 120 & 100 & ČR/EU: $\leq$25 dní/rok, 3letý průměr; WHO: Tab. 0.1 \\
O\textsubscript{3} & peak season ($\mu\text{g}\,m^{-3}$) & — & — & 60 & WHO: Tab. 0.1 (peak season) \\
O\textsubscript{3} & AOT40\hypertarget{tm:iaqlimity-a}{\hyperlink{tn:iaqlimity-a}{\tnote{a}}}\allowbreak{} ($\mu\text{g}\,m^{-3}\cdot$h) & 18000 & 18000 & — & Ochrana vegetace; 5letý průměr \\

\bottomrule
\end{tabularx}

\footnotesize \textit{Zdroj:} Vlastní zpracování na základě dat z \parencite{wolterskluwerZakon20120122012, SmerniceEvropskehoParlamentu2008, worldhealthorganizationWHOGlobalAir2021}

\begin{tablenotes}[flushleft]
  \footnotesize
  \item[a] \hypertarget{tn:iaqlimity-a}{} \textit{Pro účely tohoto zákona AOT40 znamená součet rozdílů mezi hodinovou koncentrací větší než 80 µg.m\textsuperscript{-3} (= 40 ppb) a hodnotou 80 µg.m\textsuperscript{-3} v dané periodě užitím pouze hodinových hodnot změřených každý den mezi 08:00 a 20:00 SEČ, vypočtený z hodinových hodnot v letním období (1. května - 31. července).} \parencite[příl.~1, pozn.~5]{wolterskluwerZakon20120122012}
\end{tablenotes}
\end{threeparttable}
\end{center}
\end{table}
